\subsection{River (RIV) Package}

The River Package represents a head-dependent flux between the aquifer and a
surface-water feature, limited by the riverbed bottom elevation. Its three
per-boundary inputs are the river stage, the riverbed conductance, and the
riverbed bottom elevation.

\begin{longtable}{p{0.15\textwidth} p{0.15\textwidth} p{0.60\textwidth}}
\caption{Memory-manager variables in the RIV Package that are candidates for
modification through the API. See also the variables common to all stress
packages described above.} \label{table:api-gwf-riv} \\

\hline
\textbf{Variable} & \textbf{Class} & \textbf{Guidance} \\
\hline
\endfirsthead

\hline
\textbf{Variable} & \textbf{Class} & \textbf{Guidance} \\
\hline
\endhead

\texttt{STAGE} & Period-reset &
River stage (head) for each reach, dimensioned by \texttt{MAXBOUND}. Read from
input at the start of each stress period; set it every time step to change the
stage at run time. \\

\texttt{COND} & Period-reset &
Riverbed hydraulic conductance for each reach, dimensioned by \texttt{MAXBOUND}.
Read from input at the start of each stress period; set it every time step to
change the conductance at run time. \\

\texttt{RBOT} & Period-reset &
Riverbed bottom elevation for each reach, dimensioned by \texttt{MAXBOUND}. The
exchange flux is limited when the aquifer head falls below this elevation. Read
from input at the start of each stress period; set it every time step to change
it at run time. \\

\texttt{AUXVAR} & Period-reset &
Auxiliary variable values, dimensioned by the number of auxiliary variables and
\texttt{MAXBOUND} (present when \texttt{AUXILIARY} variables are specified).
Re-read each stress period; set after \texttt{prepare\_time\_step}. Indexed by
boundary position, so a value follows the boundary index, not the cell. \\

\texttt{NODELIST} & Period-reset &
One-based reduced node number of the cell containing each boundary, dimensioned
by \texttt{MAXBOUND}. Set it after \texttt{prepare\_time\_step} to relocate a
boundary (see ``Changing Boundary Locations''). \\

\texttt{NBOUND} & Period-reset &
Number of active boundaries in the current stress period, from 0 to
\texttt{MAXBOUND}. Increase it (and fill the new entries) to activate boundaries
or decrease it to deactivate them (see ``Changing Boundary Locations''). \\

\hline
\end{longtable}

\noindent As with all of these variables, a value supplied by a time series is
refreshed each time step and overwrites API changes.
